% eksperymentalne tłumaczenie na włoski

%

\subsection{Wysokość środkowa} % lang-pl
\subsection{Altezza mediana} % lang-it

\begin{definition}
[wysokość środkowa] % lang-pl
[altezza mediana] % lang-it
Odcinek prostopadły do boku czworokąta, który kończy się na środku przeciwległego boku, nazywamy wysokością środkową. % % lang-pl
Un segmento perpendicolare a un lato del quadrilatero, che termina nel punto medio del lato opposto, si chiama altezza mediana. % % lang-it
\end{definition}

Po angielsku mamy piękne połączenie \emph{maltitude = midpoint + altitude}. % % lang-pl
In inglese abbiamo la bellissima combinazione \emph{maltitude = midpoint + altitude}. % % lang-it

\begin{definition}
% https://en.wikipedia.org/wiki/Cyclic_quadrilateral#Anticenter_and_collinearities
Niech $ABCD$ będzie czworokątem cyklicznym. % % lang-pl
Wtedy jego cztery wysokości środkowe są współpękowe, przechodzą przez punkt zwany antyśrodkiem. % % lang-pl
\index{antyśrodek}% % lang-pl
Sia $ABCD$ un quadrilatero ciclico. % % lang-it
Allora le sue quattro altezze mediane sono concorrenti, passando per un punto chiamato anticentro. % % lang-it
\index{anticentro}% % lang-it
\end{definition}

To jest \cite[zadanie 4.30]{pompe_2022}. % lang-pl
Si tratta di \cite[esercizio 4.30]{pompe_2022}. % lang-it

Joanna Jaszuńska udowodni w $\Delta_{18}^{4}$, że wysokości środkowe czworokąta (nazwie je po prostu wysokościami czworokąta) przecinają się w jednym punkcie wtedy i tylko wtedy, gdy na czworokącie można opisać okrąg; skorzysta z symetrii środkowej względem środka równoległoboku utworzonego przez środki boków, a w zadaniach pojawią się między innymi punkt Monge'a czworościanu i prosta Eulera trójkąta. % lang-pl
Joanna Jaszuńska dimostrerà in $\Delta_{18}^{4}$ che le altezze mediane di un quadrilatero (le chiamerà semplicemente altezze del quadrilatero) sono concorrenti se e solo se intorno al quadrilatero si può circoscrivere una circonferenza; si servirà della simmetria centrale rispetto al centro del parallelogramma formato dai punti medi dei lati, e tra i problemi compariranno, fra l'altro, il punto di Monge del tetraedro e la retta di Eulero del triangolo. % lang-it

\begin{proposition}
Antyśrodek pokrywa się z ortocentrum trójkąta utworzonego przez środki przekątnych oraz punkt ich przecięcia. % % lang-pl
L'anticentro coincide con l'ortocentro del triangolo formato dai punti medi delle diagonali e dal loro punto d'intersezione. % % lang-it
\end{proposition}

\begin{proposition}
Przez antyśrodek przechodzą okręgi dziewięciu punktów trójkątów $\triangle ABC$, $\triangle BCD$, $\triangle CDA$, $\triangle DAB$, co czyni go punktem Ponceleta czworokąta. % % lang-pl
Attraverso l'anticentro passano i cerchi dei nove punti dei triangoli $\triangle ABC$, $\triangle BCD$, $\triangle CDA$, $\triangle DAB$, rendendolo il punto di Poncelet del quadrilatero. % % lang-it
\end{proposition}

%